\section{Discussion and Limitations}
\label{sec:discussion}

\paragraph{What the backtest does and does not show.}
Historical validation measures whether generated questions align with
where the field actually went---strong evidence of question quality, but
conservative in one direction: a question the community never engaged
might still be valuable (the field can be collectively wrong or
resource-constrained). Conversely, engagement is partially confounded by
obviousness; a question everyone was about to ask scores well. The
premise-refuted outcome is the cleanest signal, since it requires the
system to have flagged a specific published conclusion as fragile before
the community overturned it. Scaling from 10 questions to hundreds, and
comparing against baselines (questions sampled from review-paper future
work sections; questions generated without the evidence graph), is the
necessary next step before strong claims.

\paragraph{Abstract-only extraction.}
Claims are currently extracted from abstracts. The calibration review
showed abstracts carry the load-bearing conclusions, but assumptions and
uncertainty statements often live in the body; full-text ingestion (with
the chunk-level source locations the schema already mandates) should
raise both claim coverage and adjudication quality.

\paragraph{Human adjudication is load-bearing---by design.}
Tension typing and question curation are human steps, and the reviewed
sample shows why: only 1 of 16 top candidates was a true contradiction.
The framework's contribution is to make expert judgment cheap and
auditable (typed verdicts over ranked candidates, committed as data)
rather than to eliminate it. Automating adjudication against the growing
corpus of human verdicts---with the 33\% early-precision figure as the
baseline to beat---is a measurable follow-on task.

\paragraph{LLM-judge limitations.}
The clarity dimension initially scored uniformly and was converted to a
gate; middle-band value scores compress within a few tenths of a point.
Pairwise-comparison judging and multi-judge panels are the obvious
remedies. All judge prompts and per-component scores are committed, so
these comparisons can be run retrospectively.

\paragraph{Single-domain instantiation.}
Exoplanet atmospheres was chosen because the three evidence layers align
unusually well. The framework's interfaces (claims, typed tensions,
signal-prioritized generation, two-stage ranking) are domain-agnostic,
but the topic vocabulary for tension candidates and the falsification
screens are domain-supplied; porting to a second domain is the direct
test of generality.

\paragraph{Implications for AGI research.}
The results support a specific, testable position: question-asking
competence can be decomposed into auditable stages---evidence
representation, tension detection, refinement, prioritization---each of
which can be measured and improved independently, with historical
backtesting as the end-to-end score. We suggest that benchmarks of this
form, rather than open-ended ideation judged by contemporaneous panels,
are the right instrument for tracking progress on scientific
question discovery.
